\chapter{Algèbre extérieure, déterminants}
\origpage{301--314}

Les déterminants apparaissent comme un cas particulier dans l'étude \BellaCorrection{plus générale}{plus générales} de l'algèbre extérieure, étude dont les résultats servent en calcul différentiel. C'est pourquoi je préfère les aborder de cette façon générale.

\BellaSection{1. Formes multilinéaires antisymétriques et alternées}

\medskip\noindent\textsc{Définition 9.1.} --- \emph{Soit $E$ espace vectoriel sur un corps $K$. On appelle forme $p$ linéaire sur $E$, ($p\in\mathbb N^*$) toute application $f$ de $E^p$ dans $K$ qui est linéaire par rapport à chacune de ses variables.}

Si on note $f(X_1,\ldots,X_p)$ l'image par $f$ du $p$-uplet $(X_1,\ldots,X_p)$ c'est donc que, pour chaque $i\le p$, en considérant les $X_j$ pour $j\ne i$ comme fixés, la fonction de la seule variable $X_i$ est linéaire.

\medskip\noindent\textsc{Définition 9.2.} --- \emph{Une forme $p$ linéaire, ($p\ge 2$) est dite alternée si elle s'annule lorsque deux des variables prennent des valeurs égales.}

\medskip\noindent\textsc{Définition 9.3.} --- \emph{Une forme $p$ linéaire, ($p\ge 2$) est dite antisymétrique si elle prend deux valeurs opposées lorsqu'on permute deux vecteurs.}

\medskip\noindent\textsc{Théorème 9.4.} --- \emph{Une forme alternée est antisymétrique; si le corps $K$ est de caractéristique $\ne 2$, une forme antisymétrique est alternée.}

Soit $f$ alternée, on veut permuter $X_i$ et $X_j$ avec par exemple $i<j$. En fait on calcule la valeur de $f$ en prenant en $i$ième et $j$ième position le même vecteur $X_i+X_j$. On a 0, ($f$ alternée) que l'on développe par multilinéarité en 4 termes :
\begin{align*}
f(\ldots,X_i+X_j,\ldots,X_i+X_j,\ldots)&=0\\
&=f(\ldots,X_i,\ldots,X_i,\ldots)+f(\ldots,X_i,\ldots,X_j,\ldots)\\
&\quad+f(\ldots,X_j,\ldots,X_i,\ldots)+f(\ldots,X_j,\ldots,X_j,\ldots).
\end{align*}
Les 2 termes associés aux mêmes vecteurs sont nuls, il reste donc 2 valeurs opposées quand on permute $X_i$ et $X_j$.

\emph{Réciproquement}, si $f$ est antisymétrique, et si en $i$ième position et $j$ième position, (avec par exemple $i<j$) on met le même vecteur $T$, en permutant $X_i=T$ et $X_j=T$ on obtient la même valeur, mais aussi son opposée puisque $f$ est antisymétrique.

Donc
\[
f(\ldots,T,\ldots,T,\ldots)=-f(\ldots,T,\ldots,T,\ldots)
\]
d'où
\[
2f(\ldots,T,\ldots,T,\ldots)=0.
\]
Si $\operatorname{caract}K=2$, on ne va pas plus loin, sinon on conclut $f(\ldots,T,\ldots,T,\ldots)=0$ donc $f$ est alternée. \bookend

Par la suite on supposera le corps de caractéristique différente de 2, mais on s'efforcera d'établir les résultats les plus généraux, donc sur les formes alternées.

\medskip\noindent\textsc{Théorème 9.5.} --- \emph{Soit $f$ une forme $p$ linéaire sur $E$. L'application $g$ de $E^p$ dans $K$ définie par}
\[
g(X_1,\ldots,X_p)=\sum_{\sigma\in\mathfrak S_p}\varepsilon_\sigma f(X_{\sigma(1)},\ldots,X_{\sigma(p)})
\]
\emph{est une forme $p$ linéaire alternée.}

Rappelons que $\mathfrak S_p$ est le groupe des permutations de $\{1,2,\ldots,p\}$ et que $\varepsilon_\sigma$ est la signature de la permutation $\sigma$, (voir 4.46 et 4.52).

$X_i\longmapsto f(X_{\sigma(1)},\ldots,X_{\sigma(p)})$ est linéaire car $f$ est linéaire par rapport à sa variable de rang $j=\sigma^{-1}(i)$, place où intervient $X_i$.

Donc $g$ est une combinaison de $p!$ applications linéaires en $X_i$, on a bien $g$ linéaire par rapport à chacune de ses variables.

On a $g$ alternée : supposons que $X_i=X_j=Z$, (avec $i\ne j$). Soit $T$ la transposition de $i$ et $j$ et $H=\{\mathrm{id},T\}$ le sous-groupe de $\mathfrak S_p$ engendré par $T$.

On prend l'équivalence associée à droite au sous-groupe $H$ :
\[
(\sigma\mathrel R\sigma')\Longleftrightarrow(\sigma'\sigma^{-1}\in H)
\Longleftrightarrow(\sigma'\in H\sigma)
\Longleftrightarrow(\sigma'\in\{\sigma,T\sigma\}).
\]
Chaque classe d'équivalence ayant 2 éléments, si $A$ est un ensemble de représentants de ces classes, les autres éléments sont ceux de $TA$, et les classes d'équivalences formant une partition, on a $\mathfrak S_p=A\cup TA$ d'où
\[
\begin{aligned}
g(X_1,\ldots,X_p)=\sum_{\sigma\in A}\bigl(&\varepsilon_\sigma f(X_{\sigma(1)},\ldots,X_{\sigma(p)})\\
&+\varepsilon_{T\sigma}f(X_{T\sigma(1)},\ldots,X_{T\sigma(p)})\bigr).
\end{aligned}
\]
Or $\varepsilon_{T\sigma}=\varepsilon_T\varepsilon_\sigma=-\varepsilon_\sigma$; puis si $\sigma(k)\notin\{i,j\}$, $T\sigma(k)=\sigma(k)$ et si $\sigma(k)=i$, $T\sigma(k)=j$ mais $X_{\sigma(k)}=X_i=X_j=X_{T\sigma(k)}$, alors que si $\sigma(k)=j$, $T\sigma(k)=i$ et $X_{\sigma(k)}=X_j=X_i=X_{T\sigma(k)}$.

Finalement les termes sont 2 à 2 opposés, donc $g(X_1,\ldots,X_p)=0$ lorsque $X_i=X_j$ pour $i\ne j$. \bookend

\medskip\noindent\textsc{Définition 9.6.} --- On appelle antisymétrisée de $f$, $p$ linéaire, l'application $g$ définie par
\[
g(X_1,\ldots,X_p)=\sum_{\sigma\in\mathfrak S_p}\varepsilon_\sigma f(X_{\sigma(1)},\ldots,X_{\sigma(p)}).
\]

Soient en particulier $f_1,\ldots,f_p$ des éléments du dual $E^*$, l'application
\[
f:(X_1,\ldots,X_p)\longmapsto\prod_{i=1}^pf_i(X_i)
\]
est visiblement $p$ linéaire, et son antisymétrisée s'appelle le produit extérieur des formes $f_1,\ldots,f_p$. On le note $f_1\wedge f_2\wedge\cdots\wedge f_p$.

\medskip\noindent\textsc{Définition 9.7.} --- \emph{On appelle produit extérieur de $p$ formes linéaires $f_1,\ldots,f_p$ l'application $p$ linéaire alternée $f_1\wedge\cdots\wedge f_p$ définie par :}
\[
(f_1\wedge\cdots\wedge f_p)(X_1,\ldots,X_p)=
\sum_{\sigma\in\mathfrak S_p}\varepsilon_\sigma f_1(X_{\sigma(1)})\cdots f_p(X_{\sigma(p)}).
\]

\medskip\noindent\textsc{Théorème 9.8.} --- \emph{L'application $(f_1,\ldots,f_p)\longmapsto f_1\wedge\cdots\wedge f_p$ est alternée.}

Il ne s'agit pas d'une forme, mais on a un résultat nul si deux formes $f_i$ et $f_j$ sont égales.

On suppose, avec $i\ne j$, que $f_i=f_j$ et toujours avec $T$ transposition de $i$ et $j$ et $H=\{e,T\}$ sous-groupe engendré, on considère cette fois l'équivalence associée à gauche
\[
(\sigma\mathrel{R'}\sigma')\Longleftrightarrow(\sigma^{-1}\sigma'\in H)
\Longleftrightarrow(\sigma'\in\sigma H).
\]
Si $B$ est un ensemble de représentants des classes d'équivalence, les autres éléments sont ceux de \BellaCorrectionNote{$BT$}{$\sigma B$}{Pour les classes $\sigma H=\{\sigma,\sigma T\}$, l'ensemble complémentaire des représentants est obtenu par multiplication à droite par $T$; la notation imprimée $\sigma B$ ne dépend d'aucun $\sigma$ fixé et ne peut convenir.} et $\mathfrak S_p=B\cup BT$, (partition) d'où
\[
\begin{aligned}
(f_1\wedge\cdots\wedge f_p)(X_1,\ldots,X_p)
=\sum_{\sigma\in B}\varepsilon_\sigma\bigl(&f_1(X_{\sigma(1)})\cdots f_p(X_{\sigma(p)})\\
&-f_1(X_{\sigma T(1)})\cdots f_p(X_{\sigma T(p)})\bigr).
\end{aligned}
\]
Si $k\notin\{i,j\}$, $T(k)=k$, dans les 2 produits on a $f_k(X_{\sigma(k)})$; si $k=i$, $T(k)=j$ on a d'une part $f_i(X_{\sigma(i)})$ et d'autre part $f_i(X_{\sigma(j)})$ or $f_i=f_j$ donc on trouve les facteurs $f_i(X_{\sigma(i)})$ et $f_j(X_{\sigma(j)})$.

Pour $k=j$, on retrouve $f_j(X_{\sigma(j)})$ dans le 1er produit et cette fois $f_i(X_{\sigma(i)})$ dans le 2e. Le corps $K$ étant commutatif, (hypothèse faite sur les espaces vectoriels) on a 2 produits égaux qui s'annulent, d'où $f_1\wedge\cdots\wedge f_p=0$. \bookend

\subsection*{9.9.}
On note $\bigwedge^pE^*$ l'ensemble des formes $p$ linéaires alternées sur $E$.

Cette notation se lit algèbre extérieure d'ordre $p$ sur $E$. En calcul différentiel on sera amené à définir des produits entre formes $p$ linéaires alternées et $q$ linéaires alternées d'où la dénomination extérieure.

Il est facile de voir que $\bigwedge^pE^*$ est un sous-espace vectoriel de celui des applications de $E^p$ dans $K$, et on a :

\medskip\noindent\textsc{Corollaire 9.10.} --- \emph{L'application $(f_1,\ldots,f_p)\longmapsto f_1\wedge\cdots\wedge f_p$ de \BellaCorrectionNote{$(E^*)^p$}{$E^p$}{Les arguments $f_i$ sont des formes linéaires, donc des éléments de $E^*$; le domaine correct de cette application est $(E^*)^p$.} dans $\bigwedge^pE^*$ est $p$ linéaire alternée.}

Il reste la $p$ linéarité à justifier. Or si on remplace $f_i$ par
\[\lambda_i'f_i'+\lambda_i''f_i''\]
chaque terme
\[\varepsilon_\sigma f_1(X_{\sigma(1)})\cdots f_p(X_{\sigma(p)})\]
va se scinder en deux :
\[
\lambda_i'\varepsilon_\sigma f_1(X_{\sigma(1)})\cdots f_i'(X_{\sigma(i)})\cdots f_p(X_{\sigma(p)})
\]
\[
+\lambda_i''\varepsilon_\sigma f_1(X_{\sigma(1)})\cdots f_i''(X_{\sigma(i)})\cdots f_p(X_{\sigma(p)})
\]
d'après les règles de calcul dans le corps $K$. On regroupera ensemble les termes contenant $\lambda_i'$ que l'on factorisera, ainsi que ceux contenant $\lambda_i''$ d'où finalement
\begin{align*}
f_1\wedge\cdots\wedge(\lambda_i'f_i'+\lambda_i''f_i'')\wedge\cdots\wedge f_p
&=\lambda_i'f_1\wedge\cdots\wedge f_i'\wedge\cdots\wedge f_p\\
&\quad+\lambda_i''f_1\wedge\cdots\wedge f_i''\wedge\cdots\wedge f_p.
\end{align*}

\medskip\noindent\textsc{Corollaire 9.11.} --- \emph{Si $\sigma\in\mathfrak S_p$ on a $f_{\sigma(1)}\wedge\cdots\wedge f_{\sigma(p)}=\varepsilon_\sigma f_1\wedge\cdots\wedge f_p$.}

L'application étant alternée est antisymétrique, (Théorème 9.4), et en décomposant $\sigma$ en un produit de transpositions, on obtient le corollaire, vu la définition de la signature (4.54). \bookend

\medskip\noindent\textsc{Théorème 9.12.} --- \emph{Soit $E$ espace vectoriel de dimension finie, $n$, sur le corps $K$. Si $p>n$, $\bigwedge^pE^*=\{0\}$, et si $p\le n$, $\bigwedge^pE^*$ est un espace vectoriel de dimension $\binom np$. Si $B=\{e_1,\ldots,e_n\}$ est une base de $E$, les formes $p$ linéaires alternées $e_{i_1}^*\wedge\cdots\wedge e_{i_p}^*$ associées aux $p$-uplets ordonnés \BellaCorrectionNote{$1\le i_1<i_2<\cdots<i_p\le n$}{$1\le i_1\le i_2<\cdots<i_p\le n$}{Les indices doivent être deux à deux distincts; la démonstration et la base annoncée utilisent partout les inégalités strictes.} en constituent une base.}

Si $p>n$ et si $g$ est $p$ linéaire alternée, quand on calcule $g(X_1,\ldots,X_p)$, comme la famille $\{X_1,\ldots,X_p\}$ est liée, ($p>n$, dimension de l'espace) l'un des $X_i$ est combinaison linéaire des autres.

Si
\[
X_i=\sum_{\substack{j=1\\j\ne i}}^p\lambda_jX_j,
\]
en utilisant la linéarité de $g$ par rapport à sa $i$ième variable, $g(X_1,\ldots,X_p)$ devient une somme de $p-1$ termes du type $\lambda_jg(\ldots,X_j,\ldots,X_j,\ldots)$ avec $X_j$ figurant en 2 positions, la $j$ième et la $i$ième d'où $g=0$.

On suppose $p\le n$. Il y a alors $\binom np$ façons de choisir $p$ indices distincts pris dans $\{1,\ldots,n\}$, chaque partie de $p$ indices ne s'ordonnant en croissant que d'une façon.

La famille des $e_{i_1}^*\wedge\cdots\wedge e_{i_p}^*$ avec $1\le i_1<i_2<\cdots<i_p\le n$ est libre car si on calcule
\[
(e_{i_1}^*\wedge\cdots\wedge e_{i_p}^*)(e_{j_1},\ldots,e_{j_p}),
\qquad j_1<j_2<\cdots<j_p,
\]
c'est
\[
\sum_{\sigma\in\mathfrak S_p}\varepsilon_\sigma e_{i_1}^*(e_{j_{\sigma(1)}})\cdots e_{i_p}^*(e_{j_{\sigma(p)}}).
\]
Or $e_{i_k}^*(e_{j_{\sigma(k)}})=0$ sauf si $i_k=j_{\sigma(k)}$.

Le terme générique est donc nul sauf si $i_1=j_{\sigma(1)},i_2=j_{\sigma(2)},\ldots,i_p=j_{\sigma(p)}$.

Comme les $i_k$ et $j_k$ sont indexés en croissant, ceci n'est possible que si $\sigma=\mathrm{identité}$ et $i_1=j_1,\ldots,i_p=j_p$.

D'où, pour $1\le i_1<\cdots<i_p\le n$ et $1\le j_1<j_2<\cdots<j_p\le n$ on a
\[
(e_{i_1}^*\wedge\cdots\wedge e_{i_p}^*)(e_{j_1},\ldots,e_{j_p})=
\begin{cases}
1,&i_1=j_1,\ldots,i_p=j_p,\\
0,&\text{sinon.}
\end{cases}
\]
Si on a alors une combinaison linéaire nulle
\[
\sum_{1\le i_1<i_2<\cdots<i_p\le n}\lambda_{i_1\ldots i_p}
 e_{i_1}^*\wedge\cdots\wedge e_{i_p}^*=0,
\]
en calculant en $(e_{i_1},\ldots,e_{i_p})$ il reste $\lambda_{i_1\ldots i_p}=0$ d'où la famille libre.

La famille des $e_{i_1}^*\wedge\cdots\wedge e_{i_p}^*$ est génératrice de $\bigwedge^pE^*$.

En effet soit $g\in\bigwedge^pE^*$ et considérons
\[
g'=\sum_{1\le i_1<\cdots<i_p\le n}g(e_{i_1},\ldots,e_{i_p})
 e_{i_1}^*\wedge\cdots\wedge e_{i_p}^*.
\]
La forme $g-g'$ s'annule sur chaque $(e_{j_1},\ldots,e_{j_p})$ avec $1\le j_1<\cdots<j_p$ car
\[
g'(e_{j_1},\ldots,e_{j_p})=
\sum_{1\le i_1<\cdots<i_p\le n}g(e_{i_1},\ldots,e_{i_p})
(e_{i_1}^*\wedge\cdots\wedge e_{i_p}^*)(e_{j_1},\ldots,e_{j_p})
\]
avec $(e_{i_1}^*\wedge\cdots\wedge e_{i_p}^*)(e_{j_1},\ldots,e_{j_p})=0$ sauf si $i_1=j_1,i_2=j_2,\ldots,i_p=j_p$, donc il reste $g'(e_{j_1},\ldots,e_{j_p})=g(e_{j_1},\ldots,e_{j_p})$.

Or on a le

\medskip\noindent\textsc{Lemme 9.13.} --- \emph{Un élément $h\in\bigwedge^pE^*$ est nul si et seulement si il s'annule sur tous les $(e_{j_1},\ldots,e_{j_p})$ avec $1\le j_1<\cdots<j_p\le n$.}

Il résultera du lemme que $g-g'=0$ d'où $g=g'$ sera bien combinaison linéaire des $e_{i_1}^*\wedge\cdots\wedge e_{i_p}^*$. \bookend

Justifions le lemme. Si $h=0$, il est évident que pour tout $p$-uplet
\[(e_{j_1},\ldots,e_{j_p})\]
on a $h(e_{j_1},\ldots,e_{j_p})=0$.

Réciproquement. On calcule $h(x_1,\ldots,x_p)$ avec
\[
x_k=\sum_{r_k=1}^n\xi_{r_k,k}e_{r_k}.
\]
On a
\begin{align*}
h(x_1,x_2,\ldots,x_p)
&=h\left(\sum_{r_1=1}^n\xi_{r_1,1}e_{r_1},x_2,\ldots,x_p\right)\\
&=\sum_{r_1=1}^n\xi_{r_1,1}h(e_{r_1},x_2,\ldots,x_p)
\end{align*}
par linéarité de $h$ par rapport à la première variable.

Puis $x_2=\sum_{r_2=1}^n\xi_{r_2,2}e_{r_2}$ et $h$ est linéaire par rapport à la deuxième variable. On obtient donc
\[
h(x_1,\ldots,x_p)=\sum_{r_1=1}^n\sum_{r_2=1}^n
\xi_{r_1,1}\xi_{r_2,2}h(e_{r_1},e_{r_2},x_3,\ldots,x_p)
\]
et en continuant à développer on parvient à :
\[
h(x_1,x_2,\ldots,x_p)=\sum_{r_1=1}^n\sum_{r_2=1}^n\cdots\sum_{r_p=1}^n
\xi_{r_1,1}\cdots\xi_{r_p,p}h(e_{r_1},e_{r_2},\ldots,e_{r_p}).
\]
Or chaque $h(e_{r_1},\ldots,e_{r_p})$ est nul car soit 2 des $e_{r_j}$ au moins sont égaux, et alors $h$, alternée, prend une valeur nulle; soit les indices $r_1,\ldots,r_p$ sont distincts mais il existe alors $\sigma\in\mathfrak S_p$ tel que
\[
r_{\sigma(1)}<r_{\sigma(2)}<\cdots<r_{\sigma(p)}
\]
et
\[
h(e_{r_1},\ldots,e_{r_p})=\varepsilon_\sigma h(e_{r_{\sigma(1)}},\ldots,e_{r_{\sigma(p)}})=0
\]
puisqu'ici, les indices étant ordonnés, l'hypothèse s'applique.

On a finalement $h=0$, d'où le lemme, et le théorème 9.12. \bookend

Ce théorème 9.12, dans le cas de $p=n=\dim(E)$ va nous conduire aux déterminants.

En effet, sur $E$ de dimension $n$, l'espace $\bigwedge^nE^*$ des formes $n$ linéaires alternées est de dimension 1.

Soit alors $B=\{e_1,\ldots,e_n\}$ une base de $E$, la forme $n$ linéaire alternée $e_1^*\wedge e_2^*\wedge\cdots\wedge e_n^*$ est une base de l'espace vectoriel $\bigwedge^nE^*$, (d'après le théorème 9.12).

\medskip\noindent\textsc{Définition 9.14.} --- \emph{Soit $E$ espace vectoriel de dimension finie $n$ sur le corps $K$, et $B=(e_1,\ldots,e_n)$ une base de $E$. On appelle déterminant dans la base $B$ la forme $n$-linéaire alternée $\det_B=e_1^*\wedge\cdots\wedge e_n^*$.}

Vu la définition du produit extérieur, on a donc
\[
\tag{9.15}
\det_B(X_1,\ldots,X_n)=\sum_{\sigma\in\mathfrak S_n}\varepsilon_\sigma
\xi_{1,\sigma(1)}\xi_{2,\sigma(2)}\cdots\xi_{n,\sigma(n)}
\]
car si $X_j=\sum_{i=1}^n\xi_{i,j}e_i$, on a $e_i^*(X_{\sigma(i)})=\xi_{i,\sigma(i)}$.

\BellaSection{2. Premières propriétés des déterminants}

La forme de l'expression développée de $\det_B(X_1,\ldots,X_n)$ conduit à définir le déterminant d'une matrice, puisque le second membre en fait ne dépend que des $n^2$ scalaires $(\xi_{i,j})_{\substack{1\le i\le n\\1\le j\le n}}$, c'est-à-dire de la matrice des composantes des vecteurs $X_1,\ldots,X_n$ dans la base $B$.

\medskip\noindent\textsc{Définition 9.16.} --- \emph{Soit $A=(a_{ij})_{\substack{1\le i\le n\\1\le j\le n}}$ une matrice carrée d'ordre $n$ sur un corps $K$. On appelle déterminant de la matrice $A$ le scalaire noté}
\[
\det(A)=\sum_{\sigma\in\mathfrak S_n}\varepsilon_\sigma a_{1,\sigma(1)}a_{2,\sigma(2)}\cdots a_{n,\sigma(n)}.
\]

Pour retrouver les notions du §1, on peut considérer $E=K^n$ rapporté à sa base canonique $B=(e_1,\ldots,e_n)$ et considérer les vecteurs
\[
X_j=\sum_{i=1}^na_{ij}e_i,
\]
vecteurs colonnes de $A$, on a alors :

\subsection*{9.17.}
$\det A=\det_B(X_1,\ldots,X_n)$ en fait.

\medskip\noindent\textsc{Remarque 9.18.} --- Soit $A\in M_n(K)$, $\det A$ est une somme de $n!$ termes qui sont tous, au signe près, des produits de $n$ scalaires chaque produit contenant un et un seul terme de chaque ligne, (les premiers indices varient de 1 à $n$), un et un seul terme de chaque colonne, (les seconds indices étant images de $\{1,\ldots,n\}$ par une permutation de cet ensemble). Cette remarque sert lors de l'étude du polynôme caractéristique d'un endomorphisme en dimension \BellaCorrection{finie}{fini}, (voir 10.18).

\medskip\noindent\textsc{Théorème 9.19.} --- \emph{Dans $E$ vectoriel de dimension finie, $n$, une famille de $n$ vecteurs $\{X_1,\ldots,X_n\}$ est libre si et seulement si dans une base $B$ on a $\det_B(X_1,\ldots,X_n)\ne0$ et c'est équivalent à $\det_C(X_1,\ldots,X_n)\ne0$ pour toute base $C$.}

Soit $B=(e_1,\ldots,e_n)$ une base de $E$.

Si $(X_1,\ldots,X_n)$ est une famille libre de $n$ vecteurs dans $E$ de dimension $n$, c'est une base, donc la forme $n$ linéaire alternée $X_1^*\wedge\cdots\wedge X_n^*$, associée à la base duale de $(X_1,\ldots,X_n)$, est non nulle, (elle vaut 1 sur le $n$-uplet $X_1,\ldots,X_n$) : la forme $e_1^*\wedge\cdots\wedge e_n^*$ lui est proportionnelle puisque $\bigwedge^nE^*$ est de dimension 1.

Soit $\lambda$ tel que
\[
e_1^*\wedge\cdots\wedge e_n^*=\lambda X_1^*\wedge\cdots\wedge X_n^*.
\]
On a
\[
(e_1^*\wedge\cdots\wedge e_n^*)(X_1,\ldots,X_n)=\lambda,
\]
et ce scalaire $\lambda$ est $\ne0$, puisque
\[
(e_1^*\wedge\cdots\wedge e_n^*)(e_1,\ldots,e_n)=1=
\lambda(X_1^*\wedge\cdots\wedge X_n^*)(e_1,\ldots,e_n).
\]

Si $(X_1,\ldots,X_n)$ est famille liée l'un des $X_j$ est combinaison linéaire des autres et, comme on l'a vu dans la justification du théorème 9.12, $(e_1^*\wedge\cdots\wedge e_n^*)(X_1,\ldots,X_n)=0$.

Il en résulte, par contraposée, que si $(e_1^*\wedge\cdots\wedge e_n^*)(X_1,\ldots,X_n)\ne0$, la famille $(X_1,\ldots,X_n)$ est libre. \bookend

\medskip\noindent\textsc{Corollaire 9.20.} --- \emph{Soit une matrice carrée $A$, on a $\det A=0$ si et seulement si la famille des vecteurs colonnes est liée.}

C'est encore si et seulement si l'un des vecteurs colonnes \BellaCorrection{est}{et} combinaison linéaire des autres. \bookend

Si nous avons défini le déterminant d'une matrice, c'est pour en faire quelque chose, et le relier au déterminant d'un endomorphisme en liaison avec sa matrice dans une base.

Soit donc $E$ vectoriel de dimension $n$, $u\in\operatorname{Hom}(E)$ un endomorphisme de $E$ et $g\in\bigwedge^nE^*$, (espace de dimension 1), $g$ supposée non nulle.

\subsection*{9.21.}
On définit $g_u:E^n\longmapsto K$ par :
\[
g_u(X_1,\ldots,X_n)=g(u(X_1),\ldots,u(X_n)).
\]
La linéarité de $u$ et la $n$ linéarité de $g$ impliquent la $n$ linéarité de $g_u$. De même si $X_i=X_j$, on a $u(X_i)=u(X_j)$, le 2e membre est nul, donc $g_u$ est alternée. Mais alors $g_u\in\bigwedge^nE^*$, droite vectorielle engendrée par $g$, donc il existe un facteur de proportionnalité $\alpha(u,g)\in K$ tel que
\[
g_u=\alpha(u,g)g.
\]
Le facteur ne dépend pas de $g\ne0$ dans $\bigwedge^nE^*$, car, partant de $h\ne0$ dans $\bigwedge^nE^*$, $h$ serait proportionnelle à $g$, et si $h=\lambda g$ on aura
\[
h_u(X_1,\ldots,X_n)=(\lambda g)(u(X_1),\ldots,u(X_n))=\lambda g_u(X_1,\ldots,X_n),
\]
d'où
\[
h_u=\lambda\alpha(u,g)g=\alpha(u,g)(\lambda g)=\alpha(u,g)h,
\]
le corps $K$ étant commutatif.

\medskip\noindent\textsc{Définition 9.22.} --- \emph{Soit $E\simeq K^n$ et $u\in\operatorname{Hom}(E)$. On appelle déterminant de $u$, le scalaire $\det u$ tel que, $\forall g\in\bigwedge^nE^*$, la forme $g_u$ $n$-linéaire alternée définie par $g_u(X_1,\ldots,X_n)=g(u(X_1),\ldots,u(X_n))$ vérifie $g_u=(\det u)g$.}

\medskip\noindent\textsc{Théorème 9.23.} --- \emph{Le déterminant de $u$ est le déterminant de chaque matrice de $u$ dans une base $B$ fixée de $E$.}

Car si $B=(e_1,\ldots,e_n)$ est une base de $E$, $A=(\alpha_{ij})$ est la matrice de $u$ dans cette base, les vecteurs colonnes
\[
X_j=\begin{pmatrix}\alpha_{1j}\\ \vdots\\ \alpha_{nj}\end{pmatrix}
\]
de $A$ correspondent aux vecteurs
\[
u(e_j)=\sum_{i=1}^n\alpha_{ij}e_i.
\]
Si on prend $g=e_1^*\wedge\cdots\wedge e_n^*$ pour calculer $\det u$, on a, (définition 9.17),
\[
\det A=\det_B(u(e_1),\ldots,u(e_n))=(e_1^*\wedge\cdots\wedge e_n^*)(u(e_1),\ldots,u(e_n)).
\]
Puis la relation $g_u=(\det u)g$, calculée sur le $n$-uplet $(e_1,\ldots,e_n)$ donne :
\[
g_u(e_1,\ldots,e_n)=(e_1^*\wedge\cdots\wedge e_n^*)(u(e_1),\ldots,u(e_n))
=(\det u)g(e_1,\ldots,e_n)
\]
d'où
\[
\det u=\det A.
\]
\bookend

Il résulte de ce théorème, que le déterminant de 2 matrices semblables, caractérisant le même endomorphisme dans des bases différentes, doit être le même. On va retrouver ce résultat en prouvant d'abord que :

\medskip\noindent\textsc{Théorème 9.24.} --- \emph{Soient $u$ et $v$ deux endomorphismes de $E\simeq K^n$. On a $\det(u\circ v)=(\det u)(\det v)$.}

Car, avec $g$ non nul de $\bigwedge^nE^*$, on définit $g_{u\circ v}$ par
\begin{align*}
g_{u\circ v}(X_1,\ldots,X_n)
&=g(u(v(X_1)),u(v(X_2)),\ldots,u(v(X_n)))\\
&=g_u[v(X_1),\ldots,v(X_n)]\\
&=(\det u)g(v(X_1),\ldots,v(X_n))\\
&=(\det u)g_v(X_1,\ldots,X_n)\\
&=(\det u)(\det v)g(X_1,\ldots,X_n)
\end{align*}
et ceci étant vrai quels que soient $X_1,\ldots,X_n$ c'est que $g_{u\circ v}=(\det u)(\det v)g$. Mais $g_{u\circ v}=(\det(u\circ v))g$ d'où l'égalité $\det(u\circ v)=(\det u)(\det v)$. \bookend

\medskip\noindent\textsc{Remarque 9.25.} --- Le corps $K$ étant commutatif, on a :
\[
\det(u\circ v)=\det u\det v=\det v\det u=\det(v\circ u),
\]
bien qu'en général $u\circ v\ne v\circ u$.

Soit alors $A$ et $A'=P^{-1}AP$ deux matrices semblables, on aura
\begin{align*}
\det A'&=\det(P^{-1}AP)=(\det P^{-1})(\det A)(\det P)\\
&=\det A\det(P^{-1})\det P\\
&=\det A\det(P^{-1}P).
\end{align*}
Or $P^{-1}P=I_n$ est associé à $\operatorname{id}_E$, et pour $g\in\bigwedge^nE^*$ on a $g_{\operatorname{id}_E}=g$ donc $\det(\operatorname{id}_E)=\det I_n=1$, d'où $\det A=\det A'$ résultat pressenti au théorème 9.23.

Mais au passage on vient de voir que si $P$ est inversible, on a
\[
\det(P^{-1})=\frac1{\det P}=(\det P)^{-1},
\]
avec donc $\det P\ne0$.

En fait, on a :

\medskip\noindent\textsc{Théorème 9.26.} --- \emph{Un endomorphisme $u$ de $E\simeq K^n$ est inversible si et seulement si $\det u\ne0$, et on aura $\det(u^{-1})=(\det u)^{-1}$.}

Car $u$ inversible a un inverse, $u^{-1}$, et l'égalité $u\circ u^{-1}=\operatorname{id}_E$ implique $(\det u)(\det u^{-1})=1$ d'où $\det u\ne0$ et $\det(u^{-1})=(\det u)^{-1}$.

Réciproquement, si $\det u\ne0$, avec $A$ matrice de $u$ dans une base $B=(e_1,\ldots,e_n)$ de $E$, on a
\[
(e_1^*\wedge\cdots\wedge e_n^*)(u(e_1),\ldots,u(e_n))\ne0
\]
\begin{sloppypar}
(les $u(e_j)$ associés aux vecteurs colonnes de $A$ d'après 9.17), donc la famille des vecteurs $(u(e_1),\ldots,u(e_n))$ est libre, (Théorème 9.19), c'est que l'image $u(E)$ est de dimension $n$ dans $E$ lui-même de dimension $n$, on a $\dim(\ker u)=0$ d'où $u$ injective et surjective. \bookend
\end{sloppypar}

\medskip\noindent\textsc{Corollaire 9.27.} --- \emph{Pour une matrice carrée $A$, on a $A$ inversible $\Longleftrightarrow\det A\ne0$ et $(\det A)^{-1}=\det(A^{-1})$.}

\medskip\noindent\textsc{Remarque 9.28.} --- Dans le chapitre X de la réduction des endomorphismes, on traduira, en dimension finie, la propriété $\lambda$ valeur propre de $u$, par un déterminant nul, ce qui donnera le polynôme caractéristique de l'endomorphisme $u$, (Théorème 10.18).

\medskip\noindent\textsc{Théorème 9.29.} --- \emph{Soit $E$ vectoriel de dimension finie $n$ sur le corps $K$ et $u\in\operatorname{Hom}E$, on a $\det u=\det({}^tu)$. On a de même : soit $A\in M_n(K)$, $\det A=\det({}^tA)$.}

Soit $B=(e_1,\ldots,e_n)$ une base de $E$. On considère $B^*=(e_1^*,\ldots,e_n^*)$ base duale dans $E^*$ et, pour calculer $\det({}^tu)$, si on applique ce qui précède, on doit introduire le bidual $(E^*)^*$ ainsi que la base $B^{**}$ duale de $B^*$ dans le bidual. Vous pouvez passer à la page suivante!

Si on note $\widetilde B=(\widetilde e_1,\ldots,\widetilde e_n)$ cette base $B^{**}$, $\widetilde e_j$ est définie par $\widetilde e_j(e_i^*)=0$ si $i\ne j$, 1 si $i=j$.

Soit alors $g=\widetilde e_1\wedge\cdots\wedge\widetilde e_n$ un élément de $\bigwedge^n(E^*)^*$, le scalaire $\det({}^tu)$ est tel que $g_{{}^tu}=\det({}^tu)g$, avec $g_{{}^tu}$ définie, sur $(E^*)^n$, par
\[
g_{{}^tu}(\varphi_1,\ldots,\varphi_n)=g({}^tu(\varphi_1),\ldots,{}^tu(\varphi_n))
=g(\varphi_1\circ u,\ldots,\varphi_n\circ u).
\]
En particulier, on aura
\[
g_{{}^tu}(e_1^*,\ldots,e_n^*)=\det({}^tu)\,g(e_1^*,\ldots,e_n^*)
\]
et aussi
\[
=g(e_1^*\circ u,e_2^*\circ u,\ldots,e_n^*\circ u).
\]
Or
\[
g(e_1^*,\ldots,e_n^*)=(\widetilde e_1\wedge\cdots\wedge\widetilde e_n)(e_1^*,\ldots,e_n^*)=1
\]
puisque la base des $\widetilde e_i$ est duale de celle des $e_i^*$.

Donc
\begin{align*}
\det({}^tu)&=(\widetilde e_1\wedge\cdots\wedge\widetilde e_n)(e_1^*\circ u,\ldots,e_n^*\circ u)\\
&=\sum_{\sigma\in\mathfrak S_n}\varepsilon_\sigma
\widetilde e_1(e_{\sigma(1)}^*\circ u)\cdots
\widetilde e_n(e_{\sigma(n)}^*\circ u).
\end{align*}

Soit \BellaCorrectionNote{l'isomorphisme $\theta$}{l'isomorphisme}{Le texte utilise ensuite $\theta(e_i)$ sans avoir nommé cet isomorphisme; on introduit ici la notation $\theta$.} entre le bidual et $E$ qui à un vecteur $x$ associe la forme linéaire $\widetilde x$ de $E^*\longmapsto K$ définie par
\[
\widetilde x(f)=f(x)
\]
(on vérifie facilement que $\widetilde x$ est une forme linéaire de $E^*$ dans $K$ et que $x\longmapsto\widetilde x$ est elle-même linéaire injective, car $\widetilde x(f)=0$, $\forall f\Rightarrow x=0$ sinon $x$ est alors dans une base de $E$ et la forme coordonnée suivant $x$ vaut 1 en $x$. Comme $\dim E=\dim E^*=\dim E^{**}=n$ ici, $x\longmapsto\widetilde x$ est un isomorphisme).

On a alors, pour tout $e_j^*$,
\[
\theta(e_i)(e_j^*)=e_j^*(e_i)=
\begin{cases}1,&i=j,\\0,&i\ne j,\end{cases}
\]
donc $\theta(e_i)=\widetilde e_i$ précédemment introduit.

Mais alors
\[
\widetilde e_j(e_{\sigma(j)}^*\circ u)=e_{\sigma(j)}^*(u(e_j))
\]
et
\[
\det({}^tu)=\sum_{\sigma\in\mathfrak S_n}\varepsilon_\sigma
 e_{\sigma(1)}^*(u(e_1))\cdots e_{\sigma(n)}^*(u(e_n)).
\]

Si on veut réordonner le produit pour mettre en première position $e_1^*(u(e_k))$ on doit avoir $\sigma(k)=1$ donc $k=\sigma^{-1}(1)$ d'où
\[
\det({}^tu)=\sum_{\sigma\in\mathfrak S_n}\varepsilon_\sigma
 e_1^*(u(e_{\sigma^{-1}(1)}))\cdots e_n^*(u(e_{\sigma^{-1}(n)}))
\]
et comme $\varepsilon_{\sigma^{-1}}=\varepsilon_\sigma$, et que $\sigma^{-1}$ parcourt $\mathfrak S_n$ quand $\sigma$ parcourt $\mathfrak S_n$ en posant $\sigma'=\sigma^{-1}$ on a encore
\[
\begin{aligned}
\det({}^tu)
&=\sum_{\sigma'\in\mathfrak S_n}\varepsilon_{\sigma'}
 e_1^*(u(e_{\sigma'(1)}))\cdots e_n^*(u(e_{\sigma'(n)}))\\
&=(e_1^*\wedge\cdots\wedge e_n^*)(u(e_1),\ldots,u(e_n))=\det u.
\end{aligned}
\]
Ouf! \bookend

Cette démonstration est ici ridicule par excès d'abstraction. En effet si $A$ est la matrice de $u$ dans la base $B$, on sait que ${}^tA$ est celle de ${}^tu$ dans la base duale, (voir 8.8).

On a $\det u=\det A$, (Théorème 9.23) et aussi $\det({}^tu)=\det({}^tA)$. Posons $B={}^tA$, donc $\beta_{ij}=\alpha_{ji}$ en prenant les termes généraux.

Alors
\[
\det u=\det A=\sum_{\sigma\in\mathfrak S_n}\varepsilon_\sigma
\alpha_{1,\sigma(1)}\cdots\alpha_{n,\sigma(n)},
\]
et
\[
\alpha_{1,\sigma(1)}\cdots\alpha_{n,\sigma(n)}=
\alpha_{\sigma^{-1}(1),1}\cdots\alpha_{\sigma^{-1}(n),n},
\]
en permutant les termes pour que les indices de colonnes croissent. Comme $\sigma^{-1}$ décrit $\mathfrak S_n$ lorsque $\sigma$ le décrit et que $\varepsilon_{\sigma^{-1}}=\varepsilon_\sigma$ on a
\begin{align*}
\det u&=\sum_{\sigma^{-1}\in\mathfrak S_n}\varepsilon_{\sigma^{-1}}
\alpha_{\sigma^{-1}(1),1}\cdots\alpha_{\sigma^{-1}(n),n}\\
&=\sum_{\sigma^{-1}\in\mathfrak S_n}\varepsilon_{\sigma^{-1}}
\beta_{1,\sigma^{-1}(1)}\cdots\beta_{n,\sigma^{-1}(n)}\\
&=\det B=\det({}^tu).
\end{align*}
\bookend

\medskip\noindent\textsc{Conséquence 9.30.} --- Sachons faire preuve de bons sens et utiliser les connaissances déjà acquises, parfois péniblement, pour nous faciliter le travail!

\medskip\noindent\textsc{Conséquence 9.31.} --- Ce qui a été dit pour les colonnes peut se dire pour les lignes. En particulier : une matrice $A$ est inversible, ou de déterminant non nul, si et seulement si la famille de ses vecteurs lignes est indépendante dans $K^n$.

Car $A$ inversible $\Longleftrightarrow (\det A\ne0)$, (corollaire 9.27), or $\det A=\det({}^tA)$ et $\det({}^tA)\ne0\Longleftrightarrow$ (famille des vecteurs colonnes de ${}^tA$ libre), (corollaire 9.20), et enfin les vecteurs colonnes de ${}^tA$ sont les vecteurs lignes de $A$. \bookend
